\section{Conclusion}
\label{sec:conclusion}

\subsection{Summary of Findings}

This paper has analyzed six alternative representation systems against a common
framework of five trade-off dimensions (compactness, proportionality, minority
representation, county preservation, and manipulation resistance) and identified
the constitutional availability of each.  The principal findings are:

\begin{enumerate}
  \item \textbf{No system dominates the baseline on all five dimensions.}
        Every alternative system that improves on the DIA baseline on one
        dimension degrades performance on at least one other.  The
        compactness--proportionality Pareto frontier is steep (approximately
        0.015 PP per percentage point of proportionality gained), meaning that
        even modest proportionality improvements require accepting substantially
        non-compact maps.

  \item \textbf{The Rodden sorting constraint applies to all single-member
        geographic systems.}  Changing the drawing algorithm, the objective
        function, or the jurisdictional boundaries (within single-member,
        geographic constraints) does not escape the fundamental problem that
        Democratic voters are geographically concentrated in ways that produce
        wasted votes under any compact district structure.

  \item \textbf{Multi-member districts partially escape sorting.}
        Three-member districts reduce the sorting-driven seat gap by
        approximately 40\%; five-member districts reduce it by approximately
        65\%.  This is the only mechanism available within single-election,
        geographically-bounded representation that materially improves both
        proportionality and minority representation simultaneously.

  \item \textbf{Constitutional availability constrains the reform menu.}
        National redistricting requires constitutional amendment and is
        therefore not a realistic near-term option.  Multi-member districts
        require only federal statute.  County representation in its strict
        form requires constitutional amendment.  Partisan-fairness seat
        allocation is available by federal statute under Article~I, \S4.

\end{enumerate}

\subsection{The DIA as Best Available Single-Member Option}

Given these findings, the District-Independent Algorithm --- algorithmic
single-member redistricting using edge-weighted recursive bisection without
partisan data --- achieves the best available balance of compactness, VRA
compliance, constitutional availability, and geographic stability among the
six systems examined, for systems that are:
(a) immediately available without federal legislation or constitutional amendment,
(b) compliant with all existing legal requirements (population equality,
contiguity, VRA), and
(c) resistant to partisan manipulation.

Among the six systems examined, no alternative outperforms the DIA on all four
dimensions within the same constitutional and statutory constraints.  This is
not a claim of Pareto optimality for all possible systems, but of
non-dominance within the examined portfolio.  Multi-member districts (E.1)
weakly dominate the DIA on proportionality and minority representation but
sacrifice compactness and require federal legislation to repeal 2~U.S.C.~\S2c.
This is not an argument against electoral reform; it is an argument that the
type of reform that would materially change representation outcomes
(multi-member districts) requires legislative action of a qualitatively
different kind than algorithmic adoption of a better drawing method.

\subsection{Implications for the Reform Debate}

The post-\emph{Rucho} reform debate has largely focused on independent
redistricting commissions as a remedy for partisan gerrymandering.  Our
analysis suggests a more fundamental framing: commissions adopting algorithmic
methods can eliminate drawing-process distortions, but cannot eliminate
structural distortions arising from geographic sorting.  If proportional
outcomes are the goal, the reform required is multi-member districts, not
better commissions.

This distinction is important for how reform advocates prioritize their
efforts.  Independent commissions with algorithmic methods are achievable
state-by-state through ballot initiative and are defensible against legal
challenge.  Multi-member districts require federal legislation and face
substantial political obstacles.  Both are legitimate reform goals; they are
not substitutes.

\subsection{Track E Papers}

The companion papers in Track~E develop each alternative system in full:
\citet{dellaLibera2026e1} (multi-member districts), \citet{dellaLibera2026e2}
(county representation), \citet{dellaLibera2026e3} (national redistricting),
\citet{dellaLibera2026e4} (partisan similarity), \citet{dellaLibera2026e5}
(partisan-fairness allocation), and \citet{dellaLibera2026e6} (international
applications).  The synthesis of lessons learned across all six experiments
is provided in \citet{dellaLibera2026e7} (lessons learned).
